% !TeX program = xelatex
\documentclass{jos}
\usepackage{pgfplots}
\usetikzlibrary{arrows.meta,positioning}
\pgfplotsset{compat=1.18}

% Course metadata replaces the journal publication placeholders in the template.
\makeatletter
\renewcommand{\jos@topmatter}{\noindent 天津大学 NLP 平时作业\hfill 第 8 队\par}
\renewcommand{\jos@citationbox}[2]{}
\makeatother
\fancyhead[RO]{\joskaiti 长上下文场景结构化记忆建模\quad 第 8 队}
\setlength{\headheight}{16pt}

\begin{document}
\Year{2026}
\BeginPage{1}
\CLC{TP18}
\title{长上下文场景结构化记忆建模：六篇论文阅读与比较}{结构化记忆论文阅读与比较}
\entitle{Structured Memory for Long-Context Scenarios: A Review of Six Papers}{Structured Memory Review}
\author{第 8 队（成员待填写）}{Team 8}{}
\AuthorMark{第 8 队}
\address{天津大学}{Tianjin University}
\abstract{长期交互中的记忆不仅需要保存历史，还需要保留人物、事件与时间关系，定位相关证据，并依据新信息更新状态。本文围绕这些要求，对StructMem、Hindsight、Semantic XPath、HyMEM、MetaMem和Text2Mem六篇工作进行阅读与比较，分别梳理事件组织、多通道检索、结构化访问、混合经验更新、证据利用与操作验证。比较显示，结构化记忆的价值不能仅由数据结构或单一准确率判断：来源与关系是否保真、检索范围是否正确、证据是否被合理使用以及操作效果是否符合意图，构成不同的检查层次。各论文的任务、模型、反馈条件和指标不同，原文分数不能直接汇总排名。本文以自拟的行程变更案例说明机制之间的联系与适用边界，并结合既定课设范围，说明以StructMem为主线、以其余工作辅助错误分析和局部设计的依据。相关组合仍是分析设想，本文不将其表述为已验证的新方法。}
\keywords{大语言模型；长上下文；结构化记忆；论文综述}
\maketitle

\section{背景与阅读范围}
长期对话、持续任务和跨会话助手都需要利用过去的信息。历史增长带来的困难不只在于文本长度：同一人物可能改变计划，同一事件的线索可能散落在多个会话中，系统还可能把用户陈述与自身推断混为一谈。将完整历史交给模型会增加输入开销，而将历史拆成独立片段后检索，又可能丢失人物、时间、事件归属及状态变更关系。因而，需要同时考虑信息如何进入记忆，以及后续任务如何访问和使用这些信息\upcite{xu-etal-2026-structmem,latimer-etal-2026-hindsight,liu-etal-2026-semantic}。

本文所说的结构化记忆，是对历史信息的组织单位、关系、访问范围或操作语义作出明确约束，而不限定为知识图谱。事件条目与综合摘要、带类型的记忆网络、层级树、离散策略与连续轨迹表示，以及可编辑的利用规则和操作Schema，都在不同层次提供了结构。它们的共同目的在于使后续处理有可依据的对象和边界，但任何一种表示都不能单独保证抽取正确、检索充分或最终回答可靠。

阅读范围限定为课程选题中的六篇论文：StructMem与Hindsight主要讨论长期对话记忆的组织和使用；Semantic XPath研究已有层级记忆的访问；HyMEM面向依赖截图与动作轨迹的电脑操作智能体；MetaMem优化检索后对事实的利用；Text2Mem定义记忆操作并评测执行后的状态效果\upcite{zhu-etal-2026-hybrid,xin-etal-2026-metamem,wang-etal-2026-text2mem}。这一范围覆盖了组织、访问、维护、利用和操作验证，但不是对全部智能体记忆研究的穷尽性综述。

本文按三个问题组织比较：第一，方法改变了哪一环节，需要什么输入和前提；第二，论文中的实验或消融能够支持什么判断；第三，迁入本组课设时，哪些设计可以借鉴，哪些条件尚不具备。后文各个人章节保留原论文的方法与证据，综合部分再统一讨论机制关系、共同案例和实验可比性。文中实验数字均注明来源；跨论文联系属于阅读分析，自拟案例属于解释性推演，不等同于本组已经运行的六系统对照实验。

\section{StructMem：面向长程对话的结构化记忆}
\label{sec:structmem}

StructMem\upcite{xu-etal-2026-structmem}研究长期对话中的记忆组织。它在事件内部保存事实与互动关系，并用来源时间将条目绑定；在事件之间检索历史、恢复事件语境，再生成整合记忆。细粒度条目与整合摘要共同参与问答。

\begingroup
\newcolumntype{Y}{>{\raggedright\arraybackslash}X}
\newcommand{\structmemtable}{\small\setlength{\tabcolsep}{4pt}\renewcommand{\arraystretch}{1.25}}

\subsection{研究问题与方法定位}
\subsubsection{长期对话中的关联信息}
长程对话会持续产生事实、计划、态度和互动。例如，某人提到“去年参加过一个活动”，另一人在相邻话轮中回忆了与朋友共同参加的经历。回答“两人何时一起参加活动”需要把时间表达、共同参与关系和分散的事实放在一起。仅找到与活动名称相似的句子，未必能恢复这个联系。

这类问题同时涉及两个过程。第一，系统在写入记忆时决定保留什么信息。第二，系统在读取记忆时决定给回答模型什么上下文。写入阶段丢失的关系信息，很难单靠增加检索数量补回来。StructMem 因而把记忆单元的设计作为主要切入点。
\begin{table}[!htbp]
\caption{三种记忆组织方式}
\label{tab:structmem-organization}
\centering\structmemtable
\begin{tabularx}{\linewidth}{@{}p{2cm}YYY@{}}
\toprule
组织方式 & 保存的信息 & 主要优点 & 主要代价\\\midrule
平面记忆 & 独立自然语言条目及其向量 & 构建直接，语义检索方便 & 事件上下文容易分散\\
图记忆 & 实体、关系及图结构 & 关系显式，便于结构化访问 & 实体消歧、去重与关系维护增加成本\\
StructMem & 带时间锚的事实与关系条目，加上跨事件摘要 & 保留细节，同时提供关联上下文 & 依赖抽取与整合提示的质量\\
\bottomrule
\end{tabularx}
\end{table}
% 来源：§1--3，图1，印刷页122--124
\subsubsection{结构从两个层次形成}
事件级绑定解决“同一段交流中的事实和关系怎样放在一起”。跨事件整合解决“当前交流如何与更早的经历相联系”。最终的记忆包括两类可检索内容：一类保留细节，另一类提供跨事件的解释。

\begin{figure}[!htbp]
\centering
\begin{tikzpicture}[node distance=9mm,>=Latex,every node/.style={font=\small},flowstep/.style={draw=black,align=center,minimum height=12mm,text width=3.8cm,inner sep=5pt}]
\node[flowstep] (a) {对话与来源时间};
\node[flowstep,right=8mm of a] (b) {事实条目＋关系条目};
\node[flowstep,right=8mm of b] (c) {时间锚定与事件缓存};
\node[flowstep,below=12mm of c] (d) {相关历史检索与事件重建};
\node[flowstep,left=8mm of d] (e) {跨事件整合摘要};
\node[flowstep,left=8mm of e] (f) {条目与摘要共同支持问答};
\draw[->] (a)--(b);\draw[->] (b)--(c);\draw[->] (c)--(d);\draw[->] (d)--(e);\draw[->] (e)--(f);
\end{tikzpicture}
\caption{StructMem的记忆构建与问答流程}
\label{fig:structmem-flow}
\end{figure}
% 来源：原文图2，附录A.3、A.7。

从实现上看，StructMem 属于外部记忆框架。论文的核心流程由模型调用、向量检索和数据组织组成，没有引入基础语言模型的参数训练。它仍使用向量检索，结构主要体现在事件绑定与整合内容中。

\subsection{事件级绑定：事实、关系与时间}
\subsubsection{双视角抽取}
设第 $i$ 条原话为 $m_i$，语言模型为 $\mathcal L$。事实提示 $P_{\mathrm{fact}}$ 和关系提示 $P_{\mathrm{rel}}$ 分别生成条目集合：
\begin{equation}
\Phi_i=\mathcal L(P_{\mathrm{fact}}\Vert m_i),\qquad
\Psi_i=\mathcal L(P_{\mathrm{rel}}\Vert m_i).
\end{equation}
其中，$\Phi_i$ 保存发生了什么、有什么计划或偏好；$\Psi_i$ 保存人物如何回应、支持、评价或影响对方。“双视角”对应事实内容与互动关系两个抽取目标。关系抽取关注赞扬、鼓励、感谢、兴趣、认同和回应等行为。

以研究生录取对话为例，Alice 分享自己被计算机科学项目录取的消息，Bob 祝贺她并询问研究方向。事实条目可记录录取项目与研究方向，关系条目可记录 Bob 对 Alice 的支持和兴趣。这些互动信息能帮助后续模型理解两人为何继续讨论相关计划。
\begin{table}[!htbp]
\caption{事实、关系与来源信息}
\label{tab:structmem-entries}
\centering\structmemtable
\begin{tabularx}{\linewidth}{@{}p{2.3cm}Y@{}}
\toprule
条目类型 & 条目内容\\\midrule
事实条目 & Alice 被计算机科学研究生项目录取，研究方向为机器学习。\\
关系条目 & Bob 祝贺 Alice，表达对她的信心，并询问她的研究方向。\\
来源信息 & 两类输出都携带 \texttt{source\_id}，用于对应原话及其时间。\\
\bottomrule
\end{tabularx}
\end{table}
\subsubsection{时间锚定与事件重建}
条目 $x$ 的表示包含自然语言文本、向量 $\mathbf e_x$ 和来源时间 $\tau_i$：
\begin{equation}
\mathcal M=\bigcup_i\{\langle x,\mathbf e_x,\tau_i\rangle\mid x\in\Phi_i\cup\Psi_i\}.
\end{equation}
来源时间将事实和关系绑定为一个事件单元。检索命中其中一个条目后，可以用同一时间戳找回其他条目，补全事件语境。

这里需要区分\textbf{提及时间}与\textbf{事件发生时间}。例如，2023年8月17日的一条消息说“去年参加了活动”，时间锚记录这句话的来源日期，事件年份则由相对表达推到2022年。事实抽取提示要求保留相对时间与参考日期，问答提示再要求换算为具体时间。
\subsubsection{批处理与来源追踪}
抽取过程以原话 $m_i$ 为基本单位，调用时可按话题批量处理，并按 \texttt{source\_id} 顺序输出。事实与关系分别对应两个抽取目标。来源编号、人物和时间将生成条目连接到原始对话，便于事件重建与后续核查。
% 来源：§3.1、式(1)--(2)，图4--7、图13；印刷页123--124、133--136、139

\subsection{跨事件整合与双路问答}
\subsubsection{从当前缓存找到相关历史}
系统先积累尚未整合的条目，并按来源时间排序。论文的整合触发阈值为\textbf{对话时间上的1小时}。设缓存上下文为 $C_{\mathrm{buf}}$，其文本拼接后的向量为 $q$，系统按余弦相似度检索历史中的前 $K$ 个种子条目，默认 $K=15$。

对于每个种子 $x^*$，系统找回与其共享时间戳的全部条目：
\begin{align}
E_{\tau}(x^*)&=\{x'\in\mathcal M\mid\tau(x')=\tau(x^*)\},\\
C_{\mathrm{cross}}&=C_{\mathrm{buf}}\cup\bigcup_{x^*\in S_K}E_{\tau}(x^*).
\end{align}
这一步使语义检索结果扩展为完整事件上下文。一个与问题高度相似的事实条目，可能缺少周围的赞同、回应或计划信息。时间锚能把这些条目重新聚合起来。
\subsubsection{生成有依据的整合记忆}
整合模型接收当前事件和相关历史，生成跨事件摘要。提示要求保留具体人物、日期、地点和行动计划，按时间组织内容，并在有依据时表达事件之间的联系。跨时期的引用应带有具体时间，重点放在与当前话题直接相关的历史内容上。摘要长度为200--350个英文词。

生成的摘要形成额外的抽象层，与原有事件条目共同保存。细粒度条目负责具体事实，摘要提供跨事件联系。两类记忆共同组成存储：
\[
\mathcal M_{\mathrm{all}}=\mathcal M_{\mathrm{entries}}\cup\mathcal M_{\mathrm{syntheses}}.
\]
% 来源：式(6)及问答模板；条目与整合摘要共同参与回答。
\subsubsection{回答阶段使用两路上下文}
论文在问答阶段检索\textbf{60个条目与5条整合摘要}，把它们连同问题交给回答模型。提示要求核对人物、利用时间戳换算相对日期，并给出简短答案。
\begin{table}[!htbp]
\caption{记忆整合与问答的默认配置}
\label{tab:structmem-settings}
\centering\structmemtable
\begin{tabularx}{\linewidth}{@{}p{3.2cm}p{2.2cm}Y@{}}
\toprule
参数 & 默认设置 & 所处阶段\\\midrule
时间窗口 & 1小时 & 触发记忆整合\\
历史种子数量 $K$ & 15 & 整合时检索历史，再扩展事件上下文\\
问答检索数量 & 60＋5 & 回答时读取条目与整合摘要\\
\bottomrule
\end{tabularx}
\end{table}
% 来源：§3.2、式(3)--(6)，附录 A.3，图8、图13；印刷页124、129、137、139

\subsection{案例分析：共同经历与相对时间}
\subsubsection{问题与证据}
案例问题是：Caroline 和 Melanie 什么时候一起参加过 Pride festival？关键证据分散在不同条目中，同时出现了2023年的活动记录和对“去年”活动的回忆。
\begin{table}[!htbp]
\caption{共同经历案例中的证据}
\label{tab:structmem-case-evidence}
\centering\structmemtable
\begin{tabularx}{\linewidth}{@{}p{2.2cm}Y@{}}
\toprule
证据 & 内容概述\\\midrule
事实1 & Caroline 在2023年8月11日参加了 Pride parade。\\
事实2 & Caroline 回忆自己在去年的 Pride fest 玩得很开心，该消息记录于2023年8月17日。\\
事实3 & Melanie 说自己和一群朋友在 Pride fest 度过了愉快时光，该消息同样记录于2023年8月17日。\\
互动关系 & Melanie 对 Caroline 的相关经历表示兴趣，并积极回应她参与社区活动的事情。\\
整合记忆 & 在同一叙述中连接两人的回忆，形成她们对去年共同经历的讨论。\\
\bottomrule
\end{tabularx}
\end{table}
\begin{table}[!htbp]
\caption{共同经历案例中的回答}
\label{tab:structmem-case-answers}
\centering\structmemtable
\begin{tabularx}{\linewidth}{@{}p{3cm}Y@{}}
\toprule
方法 & 模型回答\\\midrule
Flat Memory & 两人没有一起去过。\\
Graph Memory & 上个月，2023年6月。\\
StructMem & 去年，2022年8月。\\
参考答案 & 2022年。\\
\bottomrule
\end{tabularx}
\end{table}
\subsubsection{事件关联与时间推理}
问题需要同时判断“共同参加”与“参加时间”。独立事实分别描述两人的经历，互动条目补充了彼此回应的语境。整合步骤把时间相近、语义相关的信息组织成连续叙述，回答模型据此识别共同经历，并将2023年消息中的“去年”换算为2022年。

事件绑定与整合使共同经历和相对时间出现在同一回答上下文中。参考答案为2022年，StructMem给出的年份与之相符；更具体的月份仍需回溯原始对话确认。

整合记忆中的事件联系应与来源条目相连，便于回答时核对参与关系和具体日期。
% 来源：附录 A.8，表8；印刷页132

\subsection{实验设置与评价指标}
\subsubsection{数据与任务}
论文使用 LoCoMo 的长对话问答任务。数据包括10段长期对话，平均每段588轮、16,618 tokens。实验采用四类问题，共1,540题。
\begin{table}[!htbp]
\caption{LoCoMo题型分布}
\label{tab:structmem-tasks}
\centering\structmemtable
\begin{tabularx}{\linewidth}{@{}p{3.8cm}rY@{}}
\toprule
问题类型 & 数量 & 主要考察内容\\\midrule
单跳 Single-hop & 841 & 从相关记忆定位事实\\
多跳 Multi-hop & 282 & 联合多条证据得到答案\\
时间 Temporal & 321 & 时间顺序与相对日期推理\\
开放域 Open-domain & 96 & 对话记忆与开放域知识相关的回答\\\midrule
合计 & 1,540 & 单跳题约占54.6\%\\
\bottomrule
\end{tabularx}
\end{table}
\subsubsection{模型与基线}
主实验统一使用 \texttt{gpt-4o-mini} 作为基础模型，以 \texttt{text-embedding-3-small} 生成向量。比较对象覆盖原始上下文输入、RAG、平面记忆与结构化记忆。

各方法的检索配置为：MiniRAG和LightRAG各检索20条，A-Mem和LangMem各检索40条；Mem0、Mem0g、Zep和Memobase每位说话人检索10条。FullContext使用完整历史的逆时间顺序表示，表中名为OpenAI的配置直接拼接对话文本。这里的OpenAI是一个基线配置名。

不同系统使用各自的检索设置，因此主表适合比较整套系统的效果与成本。分析事件绑定和整合组件时，同框架的范式对照和消融更直接。
\subsubsection{准确率与构建效率}
准确率使用 LLM-as-a-judge：评判模型接收问题、参考答案和生成答案，输出正确或错误。评判提示允许同义表达和等价日期表述。主结果使用 \texttt{gpt-4o-mini} 判断，并用另外两种模型复核。

效率指标统计\textbf{记忆构建阶段}的输入／输出tokens、API调用次数和运行时间。它们描述把历史对话变为记忆的成本。回答阶段的tokens预算和完整服务的响应延迟，还需要单独衡量。对没有构建过程的上下文基线，以及未公开构建细节的服务，主表以横线标注相关项。
% 来源：§4.1，附录 A.2--A.5，表3，图16；印刷页125、128--130、142

\subsection{准确率结果：总分与分项表现}
表\ref{tab:structmem-accuracy}列出各方法的问答准确率，加粗表示该列最高值。
\begin{table}[!htbp]
\caption{LoCoMo问答准确率（\%）}
\label{tab:structmem-accuracy}
\centering\structmemtable
\begin{tabular}{@{}lrrrrr@{}}
\toprule
方法 & 总分 & 单跳 & 多跳 & 时间 & 开放域\\\midrule
OpenAI & 71.82 & 84.66 & 69.86 & 45.48 & 53.12\\
FullContext & 73.83 & \textbf{86.56} & 68.79 & 50.16 & 56.25\\
MiniRAG & 63.51 & 75.74 & 56.74 & 38.94 & 58.33\\
LightRAG & 68.83 & 77.53 & 66.31 & 53.89 & 50.00\\
LangMem & 58.10 & 71.12 & 62.23 & 23.43 & 47.92\\
A-Mem & 64.16 & 72.06 & 56.03 & 60.44 & 31.25\\
Mem0 & 66.88 & 72.93 & 67.13 & 59.19 & 51.15\\
MemoryOS & 58.25 & 67.06 & 56.74 & 40.19 & 45.83\\
Mem0g & 68.44 & 75.71 & 65.71 & 58.13 & 47.19\\
Zep & 75.14 & 79.79 & \textbf{74.11} & 67.71 & \textbf{66.04}\\
Memobase & 75.78 & 77.17 & 70.92 & \textbf{85.05} & 46.88\\
\textbf{StructMem} & \textbf{76.82} & 81.09 & 68.77 & 81.62 & 46.88\\
\bottomrule
\end{tabular}
\end{table}
% 来源：表1；印刷页125。加粗表示该列最高值
\subsubsection{主要结果}
StructMem的总分为76.82\%，比Memobase高1.04个百分点，比Zep高1.68个百分点，比FullContext高2.99个百分点。

分项表现更有解释力。相对FullContext，StructMem的时间推理提高31.46个百分点，而单跳下降5.47个百分点，多跳基本持平。相对Zep，StructMem时间推理更好，但多跳与开放域更低。Memobase则仍保持最高的时间推理成绩。由此可见，StructMem的收益与问题类型有关。
\subsubsection{题型差异}
全文上下文保留了大量原始细节，有利于单跳事实定位。事件结构和整合摘要帮助模型建立时间线，所以StructMem相对全文上下文的时间推理优势很突出。它在多跳任务上的机制收益，还要结合下一节的消融结果判断。

单跳题约占54.6\%，对总分影响较大。分项结果显示，StructMem的优势集中在时间推理，单跳、多跳和开放域仍有提升空间。

部分基线的汇总口径存在差异：按题型数量加权，Mem0与Zep的分类分数约为67.646\%和75.375\%，而论文报告的总分为66.88\%和75.14\%。对应行需结合有效样本和逐题结果核对。

\subsection{效率结果：构建成本与运行时间}
\subsubsection{主表中的资源消耗}
表\ref{tab:structmem-cost}统计记忆构建的资源消耗，tokens单位为百万（M），时间单位为秒。MemoryOS的总量沿用论文报告值2.868M；其输入与输出之和为2.828M，存在汇总差异。
\begin{table}[!htbp]
\caption{记忆构建的资源消耗}
\label{tab:structmem-cost}
\centering\structmemtable
\begin{tabular}{@{}lrrrrr@{}}
\toprule
方法 & 输入tokens & 输出tokens & 总tokens & 调用次数 & 时间\\\midrule
MiniRAG & 9.022 & 1.081 & 10.103 & 2,508 & 2,566\\
LightRAG & 10.014 & 1.916 & 11.931 & 13,576 & 60,469\\
LangMem & 9.873 & 1.192 & 11.066 & 5,990 & 26,281\\
A-Mem & 9.126 & 2.368 & 11.494 & 11,754 & 60,607\\
Mem0 & 10.958 & 1.239 & 12.196 & 9,181 & 30,057\\
MemoryOS & 1.889 & 0.939 & 2.868 & 5,534 & 24,220\\
Mem0g & 33.512 & 2.313 & 35.825 & 53,514 & 115,670\\
\textbf{StructMem} & \textbf{1.501} & \textbf{0.436} & \textbf{1.937} & \textbf{1,056} & 22,854\\
\bottomrule
\end{tabular}
\end{table}
% 来源：表1；印刷页125
相对Mem0g，StructMem的构建tokens减少约94.6\%，API调用减少约98.0\%，运行时间减少约80.2\%。相对Mem0，构建tokens减少约84.1\%。

运行时间具有不同的排序：MiniRAG只用2,566秒，明显短于StructMem的22,854秒。tokens数量、调用次数和耗时之间还受到模型调用方式、批处理、服务延迟与并发等因素影响。论文没有提供足够的信息，把全部耗时差异归因于其中某一因素。
\subsubsection{同一框架下的组织代价}
同框架的范式比较在LoCoMo第一段对话上拆分构建tokens：平面记忆约77.6k，StructMem约226.1k，图记忆约402.3k。StructMem的226.1k中，条目抽取约151.0k，整合约75.1k。

双视角抽取与整合增加了相对平面记忆的构建成本；相对图记忆，批量整合减少了实体、关系抽取及去重的多阶段开销。
StructMem把一部分原本留到问答时完成的关系推理，提前放进周期性的记忆构建。它用额外的组织成本换取更好的上下文，再通过批处理控制这部分成本。
% 来源：图3(a)--(b)及§4.3；印刷页126

\subsection{消融与参数分析}
\subsubsection{事件绑定与跨事件整合的贡献}
\begin{table}[!htbp]
\caption{记忆组织消融：准确率（\%）及百分点变化}
\label{tab:structmem-ablation}
\centering\structmemtable
\begin{tabular}{@{}lrrrr@{}}
\toprule
配置 & 单跳 & 多跳 & 时间 & 开放域\\\midrule
Flat Memory & 78.83 & 66.31 & 78.50 & 46.88\\
Graph Memory & 80.50 & 66.67 & 76.64 & 48.96\\
去掉跨事件整合 & 80.86 & 66.31 & 79.44 & 46.88\\
StructMem & 81.09 & 68.77 & 81.62 & 46.88\\\midrule
加入跨事件整合的变化 & +0.23 & +2.46 & +2.18 & 0.00\\
\bottomrule
\end{tabular}
\end{table}

保留事件绑定、去掉跨事件整合后，单跳与时间任务仍高于Flat Memory。进一步加入整合，多跳提高2.46个百分点，时间提高2.18个百分点，单跳只增加0.23个百分点。这与方法设计相符：跨事件摘要主要帮助联合分散信息。去掉与加入整合的两种配置中，开放域分数均为46.88\%。
\subsubsection{历史种子数量与效果}
\begin{figure}[!htbp]
\centering
\begin{tikzpicture}
\begin{axis}[width=.85\linewidth,height=4.5cm,xlabel={整合阶段的历史种子数量 $K$},ylabel={准确率（\%）},xmin=0,xmax=20,ymin=75.4,ymax=77.05,xtick={0,5,10,15,20},ytick={75.5,76,76.5,77},axis lines=left,grid=major,grid style={gray!15},tick label style={font=\small},label style={font=\small}]
\addplot[color=black,mark=*,thick,nodes near coords,point meta=y,every node near coord/.append style={font=\scriptsize,anchor=south},/pgf/number format/fixed,/pgf/number format/precision=2] coordinates {(0,75.71)(5,76.43)(10,76.49)(15,76.82)(20,76.23)};
\end{axis}
\end{tikzpicture}
\caption{历史种子数量对问答准确率的影响（纵轴范围75.4--77.05）}
\label{fig:structmem-seeds}
\end{figure}
% 来源：原文图3(d)。

$K=0$时仍整合当前缓存，只是不引入历史事件。$K=15$达到图中最高的76.82\%，$K=20$则回落到76.23\%。这表明更多历史上下文不一定继续改善效果，选择范围需要控制。

在检索规模分析中，平面检索从10条增加到60条，准确率由70.91\%提高到75.78\%；继续增加到100条时为75.71\%，相关tokens量则从3.65M增到5.62M。该结果支持适度增加检索覆盖率，也说明覆盖率之外仍存在信息组织问题。
% 来源：表2，图3(c)--(d)；印刷页125--126

\subsection{生成记忆的忠实性与评判稳定性}
\subsubsection{抽取和整合分别接受检验}
新的记忆由语言模型生成，因此需要检查生成内容是否有原始证据。忠实性评测分成两步：先检查单个条目是否忠于原话，再检查跨事件摘要是否引入了错误联系。

在事件抽取阶段，三个评判模型报告的幻觉率分别为2.63\%（GPT）、1.61\%（Qwen）和2.84\%（DeepSeek），原文汇总的平均值为2.36\%。这描述的是抽取条目的忠实性，与问答错误率是不同指标。

在整合阶段，评判模型比较“只整合当前缓存”与“加入15个历史种子”的摘要，识别新增历史联系。作者进一步删除提示中的具体时间引用与具体联系要求，得到无约束版本。
\begin{table}[!htbp]
\caption{来源约束对错误关联的影响}
\label{tab:structmem-faithfulness}
\centering\structmemtable
\begin{tabular}{@{}lrrr@{}}
\toprule
评判模型 & 有约束的错误联系 & 无约束的错误联系 & 错误率变化\\\midrule
GPT & 6/982（0.61\%） & 72/966（7.45\%） & +6.84个百分点\\
Qwen & 33/968（3.41\%） & 213/1065（20.00\%） & +16.59个百分点\\
DeepSeek & 13/358（3.63\%） & 104/687（15.14\%） & +11.51个百分点\\
\bottomrule
\end{tabular}
\end{table}
这组实验支持具体时间与来源关联约束的作用。不同评判模型识别出的联系数量不同，因此表中保留各自的分子和分母，分别比较约束前后变化。
\subsubsection{更换问答评判模型后的排名}
\begin{table}[!htbp]
\caption{不同评判模型下的准确率（\%）}
\label{tab:structmem-judges}
\centering\structmemtable
\begin{tabular}{@{}lrrr@{}}
\toprule
评判模型 & StructMem & Memobase & 分差（百分点）\\\midrule
gpt-4o-mini & 76.82 & 75.78 & +1.04\\
Qwen2.5-32B-Instruct & 77.01 & 76.36 & +0.65\\
DeepSeek-V3.2 & 79.35 & 79.29 & +0.06\\
\bottomrule
\end{tabular}
\end{table}
三个评判模型下，StructMem的总分都略高于Memobase。作者报告Fleiss' $\kappa=0.8341$，表明这些模型的判断具有较高一致性。与此同时，DeepSeek评判下的分差只有0.06个百分点。排名相近的系统还需要逐题分析和重复实验来理解差异来源。
% 来源：附录 A.5--A.7，表4--7，图17--20；印刷页129--131、143--146

\subsection{方法局限与课设启发}
\subsubsection{关系保存与多粒度记忆}
StructMem将关系信息的保存提前到记忆形成阶段。人物的支持、回应和共同回忆经常出现在话语互动中；保留这些关系有助于重建事件语境，再通过整合连接相关历史。

原子条目和整合摘要提供了互补粒度：条目保存具体细节，摘要组织分散信息。两种粒度共同检索，兼顾事实定位与跨事件理解，也提供了清晰的消融维度。
\subsubsection{更新、追溯与泛化}
\textbf{旧记忆如何修订。}作者在局限部分指出，当前框架缺少显式的冲突消解与记忆更新机制。用户更改计划或偏好后，旧摘要可能继续被检索。问答提示中的“优先最近记忆”提供了回答规则，存储层仍需要版本、失效标记或重新整合机制。

\textbf{摘要联系如何追溯。}时间锚和提示约束降低了无依据联系，但摘要中的每项关联若能对应到具体来源条目，检查和更新会更加直接。尤其当某条历史事实失效时，系统应能找到受它影响的摘要。

\textbf{收益能否迁移。}论文主要在LoCoMo上测试，基础模型固定为gpt-4o-mini。更多任务类型、不同模型和更长历史能够帮助判断方法适用范围。现有消融尚未逐项拆开“关系抽取”“时间锚定”“摘要检索”的贡献，后续可以在固定预算下继续比较。
\subsubsection{适合课设的小规模验证}
以同一基础模型和同一份数据，比较事实记忆、事件记忆、事件记忆加跨事件整合三种配置。先固定检索预算，再分别去掉关系条目、时间信息和整合摘要。记录四类准确率、构建与问答tokens、调用次数，并保留失败题的检索上下文。

还可加入计划变更对话，观察旧事实和旧摘要的影响，并比较来源追踪与摘要失效策略。
% 来源：§3--4，Limitations，附录 A.3、A.7；印刷页123--127、129、139

\subsection{小结}

StructMem通过事件内绑定保留人物互动与事实语境，通过跨事件整合建立时间联系。主结果、组件消融和忠实性检查分别展示了整体表现、整合模块的作用与生成关联的可靠性。

该方法适合分析时间关系和跨事件证据组织。进一步研究可聚焦关系抽取、时间锚和摘要检索的独立贡献，以及计划变更后旧条目与摘要的一致性维护。
\endgroup

\section{Hindsight：保留、召回与反思的结构化智能体记忆}
\label{sec:hindsight}

Hindsight\upcite{latimer-etal-2026-hindsight}面向跨多轮对话和多个会话保持记忆的智能体。它把记忆组织为四类网络，并用retain、recall、reflect三种操作分别处理信息写入、证据检索和基于记忆的回答。

\subsection{研究问题与四类记忆}

作者将现有对话记忆的困难概括为三个相连的问题：长历史中怎样保留有用信息，查询时怎样跨会话找回相关信息，以及怎样在回答时避免只凭局部线索而产生全局不一致。单纯把片段放进向量库并取回 top-$k$，还可能混淆用户说过的事实与智能体形成的判断，也难以表达时间和实体之间的联系。

Hindsight 在每个 memory bank 中维护四类逻辑网络：
\begin{table}[!htbp]
\centering
\caption{Hindsight的四类记忆}
\label{tab:hindsight-types}
\begin{tabularx}{\linewidth}{@{}p{2.5cm}XX@{}}
\toprule
网络 & 保存内容 & 例子或作用 \\
\midrule
World & 关于外部环境的客观事实 & Alice 在 Google 的 Mountain View 团队工作。 \\
Experience & 智能体自身采取过的行动 & 智能体曾向 Alice 推荐徒步旅行。 \\
Observation & 根据多条事实整合出的实体摘要 & Alice 是专长于机器学习的软件工程师。 \\
Opinion & 带置信度、可随证据变化的主观判断 & “Python 适合数据科学”，并记录置信度。 \\
\bottomrule
\end{tabularx}
\end{table}
World与Experience记录事实或经历，Observation汇总实体证据，Opinion保存主观判断。这些类型区分了事实、概括与观点，便于检查和更新。记忆单元带有事件发生时间区间和系统获知时间，并通过时间、语义、实体和因果关系连接成图。

\subsection{三种操作：写入、检索与生成}

\textbf{Retain} 接收对话或文档，由模型抽取较粗粒度的叙事事实，再进行分类、实体识别与消歧、时间标注和关系建图。作者称每段对话通常抽取 2--5 条事实，以减少记忆切分方式对后续检索的影响。这里记录的时间至少要区分事件发生时间与智能体提及或获知该信息的时间。

\textbf{Recall} 接收查询和 token 预算，返回预算内排序后的记忆。它并行使用四类检索：HNSW 向量相似度、BM25 关键词匹配、沿实体／时间／因果连接进行图遍历，以及按事件时间范围过滤。不同通道的排序由 Reciprocal Rank Fusion 合并，再由 cross-encoder 重排。图遍历通常限制在 2--3 跳；时间解析先处理常见表达，规则无法处理时再调用轻量 T5 模型。token 预算让调用者限定送入回答模型的上下文规模。

\textbf{Reflect} 将查询、取回的记忆和可配置的行为档案交给生成流程。行为档案由怀疑程度、字面倾向和共情程度三个 1--5 的参数组成。Opinion 与事实分开存储，置信度在 $[0,1]$ 范围内；新证据支持某个判断时提高置信度，形成反证时降低置信度。后台整合还会保存支持性引文、证据数量和从 new 到 stale 的新鲜度趋势，使判断可以回溯到证据。

论文把 retain 和 recall 放在 TEMPR（Temporal Entity Memory Priming Retrieval）组件中，把 reflect 放在 CARA（Coherent Adaptive Reasoning Agents）组件中。端到端流程可概括为“对话进入四类记忆网络—多通道召回—依据证据生成并更新意见”。作者使用 PostgreSQL、pgvector 与全文检索索引作为存储基础，并报告：包含 10,000 条记忆的 bank，其 recall 检索延迟低于 200 ms，但该数字不包含回答模型的推理时间。

\subsection{实验设置与结果}

论文在LongMemEval与LoCoMo上评估长期对话记忆。LongMemEval的S设置包含500道问题，对话最长约150万tokens；LoCoMo使用多会话对话，每段最多35个会话。表\ref{tab:hindsight-results}列出不同模型配置下的回答准确率。

\begin{table}[!htbp]
\centering
\caption{Hindsight在两项基准上的回答准确率（\%）}
\label{tab:hindsight-results}
\begin{tabular}{lrrr}
\toprule
基准 & Hindsight 20B & Hindsight 120B & Hindsight Gemini-3 \\
\midrule
LongMemEval & 83.6 & 89.0 & 91.4 \\
LoCoMo & 83.2 & 85.7 & 89.6 \\
\bottomrule
\end{tabular}
\end{table}

在LongMemEval上，20B模型配合Hindsight达到83.6\%，120B模型为89.0\%，Gemini-3为91.4\%；在LoCoMo上，相应结果为83.2\%、85.7\%和89.6\%。20B配置在LongMemEval时间推理和多会话题上均为79.7\%，偏好题为66.7\%，反映出不同题型仍有差异。

作为参照，全文GPT-4o在LongMemEval上的分数为60.2\%，Supermemory-Gemini-3为85.2\%；LoCoMo中Zep、Memobase和Backboard分别为75.1\%、75.8\%和90.0\%。这些结果涉及不同基础模型和系统配置，适合描述各配置的表现；记忆机制的独立贡献则需在固定模型与预算下比较。

\subsection{方法局限与课设启发}

论文明确列出的局限包括：事实抽取、实体消歧和意见更新依赖语言模型，错误可能沿记忆图传播；评测只覆盖英文基准；Opinion 的主观质量和置信度校准还没有通过正式用户研究验证；时间解析可能漏掉含混或文化相关的表达。持续保存个人偏好与意见也带来隐私和偏见风险，论文指出系统没有内置 PII 检测或自动脱敏机制。

评测材料也有待完善：系统演示论文引用的完整评估条目仍带有arXiv编号占位符，且正文与附录中的部分LongMemEval分类基线数值不一致。复核组件效果时，需要补充完整实验配置和逐题结果。

Hindsight对课程设计的启发集中在记忆类型和检索组织：分别保存事实、经历、实体摘要与观点，利用时间和实体关系补充向量检索，并通过token预算控制回答上下文。可在固定模型、数据和预算下逐项比较检索通道，记录准确率、无关条目比例、延迟与token开销。

% 原文定位：PDF 第1--3页，印刷页275--277；PDF 第3--4页，§3.1、图1；PDF 第3--4页，§3.2--3.3；PDF 第3、5页，§3.2、§4.1；PDF 第4--5页，图1、§4.1--4.2；PDF 第7页及第11页，§6.1--6.3、表4；PDF 第8页，印刷页282；PDF 第7--9、11页，§6.2、表4及参考文献；PDF 第4页；PDF 第6页；PDF 第11页。

\section{Semantic XPath}
Semantic XPath\upcite{liu-etal-2026-semantic}研究结构化对话记忆的访问。它结合路径、位置条件和连续语义评分，从已有树形记忆中选择相关节点及子树，支持下游回答和状态更新。

\subsection{研究问题与方法定位}

长期任务导向对话需要反复查看、修改和回溯同一个任务对象。旅行计划中的一项活动不仅包含文本，还属于特定日期、行程与历史版本。全量上下文保留信息较多，但输入量随历史增长；扁平检索降低输入量，却可能取回语义相似而所属范围错误的条目。比如“在会议安排最密集的一天增加咖啡休息”，需要比较各天的活动，而不仅是找一条包含会议内容的文本。

查询负责定位访问范围，插入、删除及版本维护由下游步骤完成。

\subsection{树形数据模型与查询语法}

记忆表示为有根树 $T=(V,E,r)$。非根节点只有一个父节点，节点具有类型和文本属性；模式 $\Sigma$ 由记忆数据已有的节点类型、属性和父子关系决定，而非要求人工预先设计一个通用模式。旅行示例采用
\[
\text{Itinerary}\rightarrow\text{Version}\rightarrow\text{Day}\rightarrow\text{POI},
\]
分别表示计划、版本、日期和活动；层级主要是组成关系，不只是对下层文本做有损摘要。

查询是步骤序列，每一步写为 $S=A\,N\,[R]\,[P]$：路径与类型必需，位置选择器和语义条件可选。
\begin{table}[H]
\centering
\caption{Semantic XPath 查询步骤的组成}
\label{tab:sxpath-grammar}
\begin{tabularx}{\linewidth}{@{}p{20mm}p{43mm}X@{}}
\toprule
组成 & 示例写法 & 作用 \\
\midrule
路径 $A$ & \texttt{/}、\texttt{//} & 访问直接子节点或全部后代 \\
类型 $N$ & \texttt{Day}、\texttt{POI}、\texttt{*} & 按节点类型筛选或匹配任意类型 \\
位置 $R$ & \texttt{[i]}、\texttt{[-i]}、\texttt{[i:j]} & 按顺序、倒序或区间选取候选 \\
语义 $P$ & \texttt{[node\textasciitilde="workshop"]} & 计算候选与语义条件的相关程度 \\
\bottomrule
\end{tabularx}
\end{table}

例如，\texttt{//Day[3]/POI[node\textasciitilde="workshop"]}先选第三个匹配日期，再对该日期下的活动评分。存在多个行程或版本时，查询需先限定目标范围。语义条件$1-P$对原条件评分取反，提高低分候选的优先级。

\subsection{语义评分与带权执行}

语义相关度定义为 $\operatorname{Rel}(u,P)\in[0,1]$。局部评分依据一个属性或完整节点表示；聚合评分先在当前节点下执行内部结构查询，再用平均、最小、最大或几何平均聚合证据；多个条件还可用乘积、平均及取反组合。

\begin{quote}\small\ttfamily
//Day[avg(/POI[node\textasciitilde="conference"])]
\end{quote}
这个查询根据活动证据评价日期。示例中第二天的三个活动分数为$0.603$、$0.482$、$0.608$，平均约$0.565$，在候选日期中得分最高。该值衡量活动与会议主题的平均相关度。

执行器从根节点及权重1出发，维护带权集合 $W\subseteq V\times[0,1]$，按步骤进行路径扩展、类型筛选、位置选择与语义评价。语义步骤更新为
\begin{equation}
w_{\mathrm{new}}=w_{\mathrm{old}}\operatorname{Rel}(u,P).
\end{equation}
聚合条件递归访问下层证据。阈值、返回数量和空集处理由具体任务设置，查询质量则可从目标范围、节点选择和证据完整性三个方面检查。

\subsection{范围约束示例}

以行程调整为例，假设当前版本的第2天安排海滨公园，第5天安排山地徒步，第7天安排植物园和海港游船。用户要求找出“第7天下雨时需要调整的户外活动”。纯语义排序可能混入其他日期的活动；结构化访问先限定第7天，再评价活动语义，目标为植物园与游船。

简单日期约束也可由带元数据过滤的扁平RAG完成。比较查询机制时，可将其作为基线，进一步考察嵌套范围、下层证据聚合等结构访问需求。

\subsection{实验设置与指标口径}

作者评估旅行行程、待办清单和家庭配餐三个自建场景，每场景20个单轮对话、5个多轮对话；另在 LoCoMo 中使用原生的 Memory--Conversation--Session--Turn 层级。回答模型为 GPT-5 mini 与 Gemini 3 Flash。相似度评分器（Sem）和 Flat RAG 使用 Qwen3-Embedding-8B；蕴含评分器（Entail）使用 \texttt{facebook/bart-large-mnli}。

自建场景使用任务通过率，要求检索到正确结构化数据且输出满足任务约束；LoCoMo使用模型裁判评定的回答得分。LoCoMo的平均输入tokens包含XPath查询生成输入，自建场景则报告平均token用量。成本比较还需计入嵌入、蕴含评分和查询执行，并在同一模型内对齐分词与统计口径。

\begin{table}[H]
\centering
\caption{三个自建场景的平均单轮结果}
\label{tab:sxpath-single}
\begin{tabular}{@{}llrr@{}}
\toprule
回答模型 & 方法 & 通过率（\%） & 平均tokens \\
\midrule
GPT-5 mini & In-context & 70.0 & 26\,216 \\
 & Flat RAG & 30.0 & 3\,250 \\
 & XPath Sem & 66.3 & 6\,166 \\
 & XPath Entail & 83.0 & 5\,691 \\
\midrule
Gemini 3 Flash & In-context & 79.0 & 25\,387 \\
 & Flat RAG & 37.7 & 3\,361 \\
 & XPath Sem & 64.3 & 6\,043 \\
 & XPath Entail & 70.7 & 5\,830 \\
\bottomrule
\end{tabular}
\end{table}

单轮GPT-5 mini配置中，Entail相对Flat RAG提高53.0个百分点，相对增幅约176.7\%；token用量约为全量上下文的21.7\%。Gemini配置中，全量上下文的79.0\%高于Entail的70.7\%，评分机制的收益随模型变化。

\begin{table}[H]
\centering
\caption{三个自建场景的平均多轮结果}
\label{tab:sxpath-multi}
\begin{tabular}{@{}llrr@{}}
\toprule
回答模型 & 方法 & 通过率（\%） & 平均tokens \\
\midrule
GPT-5 mini & In-context & 80.0 & 69\,097 \\
 & Flat RAG & 26.7 & 3\,327 \\
 & XPath Sem & 70.0 & 6\,294 \\
 & XPath Entail & 83.3 & 5\,850 \\
\midrule
Gemini 3 Flash & In-context & 56.7 & 76\,603 \\
 & Flat RAG & 26.7 & 3\,240 \\
 & XPath Sem & 76.7 & 6\,092 \\
 & XPath Entail & 80.0 & 5\,993 \\
\bottomrule
\end{tabular}
\end{table}

多轮Entail在两种模型下取得最高平均通过率，token用量约为全量上下文的8.5\%与7.8\%。Sem与Entail的差异说明评分器选择会影响任务效果。单轮与多轮实验使用不同请求，分别反映两类任务设置。

\begin{table}[H]
\centering
\caption{LoCoMo结果（XPath输入含查询生成）}
\label{tab:sxpath-locomo}
\begin{tabular}{@{}llrr@{}}
\toprule
回答模型 & 方法 & 回答得分 & 平均输入tokens \\
\midrule
GPT-5 mini & In-context & 74.55 & 24\,857 \\
 & Flat RAG & 35.71 & 1\,341 \\
 & Semantic XPath & 65.19 & 3\,032 \\
\midrule
Gemini 3 Flash & In-context & 87.01 & 27\,499 \\
 & Flat RAG & 41.49 & 1\,996 \\
 & Semantic XPath & 73.25 & 3\,870 \\
\bottomrule
\end{tabular}
\end{table}

LoCoMo上XPath的总体得分处于两种基线之间，输入分别约为全量上下文的12.2\%与14.1\%。这说明已有会话层级可用于访问，但准确性仍有损失。例如 GPT-5 mini 多跳题得分为43.26，低于全量上下文58.87，证据整合仍有差距。作者的失败分析将错误分为查询生成、检索和回答生成，认为检索是主要剩余来源，包括无命中、证据不全及选错节点或发言。

\subsection{方法局限与课设启发}

结构范围与连续语义评分支持不同粒度的访问，路径与评分轨迹便于解释结果。方法依赖可靠层级，跨树关系、模糊版本和复杂多跳仍有困难。现有评测规模较小，重复实验、置信区间及本地评分开销也有待补充。

StructMem侧重写入时保存事件关联，Semantic XPath侧重读取时限定范围。课程设计可明确日期与版本约束，保留节点及来源标识，分别检查范围、证据完整性和回答错误，并以元数据过滤基线评估层级访问的贡献。

% 原文定位：PDF第1--2页，图1--2；PDF第3页，\S3.1；PDF第3--4页，\S3.2；PDF第10页；PDF第4页，\S3.3、图2；PDF第5页，\S4.1；PDF第6--7页，\S4.3、图8。

\section{HyMEM：电脑操作智能体的混合自演化记忆}
\label{sec:hymem}

HyMEM\upcite{zhu-etal-2026-hybrid}面向需要观察截图并连续操作网页的电脑操作智能体。历史经验既包含可复用策略，也包含界面与动作细节；任务推进后，最初检索的经验可能失效。系统因此结合离散策略与连续轨迹表示，分别更新长期记忆和当前任务的工作记忆。

\subsection{问题背景与混合表示}

输入包括任务描述、当前截图及历史成功操作轨迹，输出为下一步可执行操作。例如，搜索商品和结账所需经验不同，但可能需要保留同一预算约束。纯文本总结容易丢失视觉细节，单纯连续表示则缺少清晰的策略组织和可解释索引；一次性检索也难以覆盖交互阶段变化。这些问题构成混合表示与两种更新机制的动机。

论文在图 $G=(V,E)$ 中组织经验，将一条轨迹的逻辑表示写为
\begin{equation}
v_i=(c_i,A_i,m_i),
\end{equation}
其中 $c_i$ 为高层策略，$A_i$ 为动作、界面元素或领域标签等属性集合，$m_i$ 为连续轨迹表示。共享相同属性的轨迹间建立无向边，并通过策略、属性与轨迹的关联支持跨经验检索。策略与属性提供语义锚点；借鉴CoMEM的连续编码保留难以用文本概括的视觉及动作细节，每条轨迹压缩为8个嵌入向量并接入智能体输入。

检索使用独立的索引向量：归一化的任务文本CLIP表示与起始截图CLIP表示拼接后，通过FAISS寻找相似经验。连续记忆$m_i$则负责向智能体提供检索到的轨迹细节。

\subsection{全局记忆演化与结构化检索}

更新长期图时，新轨迹先检索相关经验，再由视觉语言模型判断信息增量。检索提供候选集合，模型进一步判断经验是否重复。三种更新为：
\begin{itemize}
    \item \textbf{ADD}：出现新策略或属性，增加对应节点与连接。
    \item \textbf{MERGE}：已有策略获得互补证据，如界面变体或另一条操作路径，将证据并入相应策略及属性关联。
    \item \textbf{REPLACE}：在同等条件下，新轨迹以更少步骤、更高成功率或更低风险实现相同策略，用其替换过时证据并刷新关联。
\end{itemize}
这些操作在长期图中积累可迁移经验并减少冗余。更新效果取决于模型对信息增量和轨迹质量的判断。

推理时，系统依据任务和截图寻找相似种子，再沿图扩展邻居并按相似度排序。主实验每次使用10条轨迹，包括5条语义相似种子和从其1跳邻居中选出的5条轨迹。一般算法允许进一步扩展邻居。

\subsection{混合工作记忆及局部刷新}

检索后形成两个互补视图：策略与属性归纳为指导指令（Guidance Instructions），作为文本写入系统提示；轨迹连续嵌入则与模型输入连接。前者帮助规划，后者提供细粒度视觉和动作信息，二者构成当前任务的工作记忆。

每个动作后，视觉语言模型比较前后观测与当前指导，判断是否发生阶段转换。若任务由搜索转入结账，先保留长期目标和约束，丢弃不相关内容，再按新状态检索，合并指导并刷新连续嵌入。这是\textbf{任务内部的上下文更新}，与长期图的ADD/MERGE/REPLACE属于不同层次。

\subsection{实验设置与结果}

实验覆盖WebVoyager、Multimodal-Mind2Web和MMInA三个多模态网页智能体基准。前两者根据任务描述与截图轨迹进行模型判分；MMInA比较最终回答和标准答案。主指标是任务成功率，Overall按可用领域取平均。各方法在统一环境、抽样和评价方式下实现。

记忆来源为GUIAct、Mind2Web训练集及作者智能体生成的成功轨迹，共2883条；主实验图含1858个节点和超过100万条边。执行采用ReAct式接口，包含SOM界面标记及UI-INS-7B定位回退。训练更新记忆编码器的Q-Former和LoRA层，可训练参数约占总参数的1.2\%，通常训练1个epoch。系统效果因此同时依赖记忆数据、编码器训练和界面定位能力。

\begin{table}[!htbp]
\centering
\caption{不同记忆表示下的Overall任务成功率（\%）}
\label{tab:hymem-results}
\begin{tabular}{lrrrr}
\toprule
基础模型 & 无记忆 & 离散记忆 & 连续记忆 & HyMEM \\
\midrule
Qwen2.5-VL-7B & 12.5 & 17.0 & 21.7 & 35.0 \\
Qwen3-VL-8B & 21.9 & 21.3 & 23.3 & 33.8 \\
UI-TARS-1.5-7B & 17.1 & 15.4 & 28.7 & 31.2 \\
\bottomrule
\end{tabular}
\end{table}

Qwen2.5-VL-7B由12.5\%升至35.0\%，提高22.5个百分点，混合记忆高于单独的离散或连续记忆。不同模型和领域的收益有差异：Qwen3-VL-8B离散记忆的21.3\%略低于无记忆的21.9\%，部分领域采用HyMEM后也未改善。

在Amazon与Google Maps的更新分析中，全局演化使成功率分别由48.8\%升至61.0\%、43.9\%升至51.2\%；局部刷新分别达到56.1\%和48.8\%。两项分析各自使用静态记忆作为参照，展示跨任务更新和任务内刷新的作用。

\subsection{方法局限与课设启发}

HyMEM结合可解释策略与连续视觉细节，并通过长期图更新和任务内刷新适应状态变化。其开销包括记忆构建、编码器训练、图扩展和执行期间的模型判断。图关联与更新依赖标签质量及模型判断，错误记忆传播、阶段误判和跨网站迁移是进一步评测的重点。

对文本事件记忆而言，可借鉴HyMEM对语义种子与结构邻居、长期更新与工作记忆刷新的区分。课程设计可在固定检索预算下比较邻居扩展带来的证据覆盖、无关信息和调用成本；连续轨迹编码则适用于具有截图和动作序列的多模态任务。

% 原文定位：PDF第1--3页，表1；PDF第3--5页，\S3.1--3.3、图1；PDF第3页，式(1)；PDF第4页，图1与Memory Update；PDF第4--5页，Structured Retrieval及Memory Settings；PDF第5页，On-the-fly Working Memory Refresh；PDF第5--7页，\S4.1--4.2与表2；PDF第5页，\S4.1；PDF第7页，\S4.3.1、图2。

\section{MetaMem：由事实存储转向事实利用}
\label{sec:metamem}

MetaMem\upcite{xin-etal-2026-metamem}研究检索后的事实利用：模型即使找到相关记忆，仍可能误解时间、数量或从属关系。该方法将可复用的知识使用经验保存为元记忆，并通过答案反馈迭代优化。

\subsection{研究问题与两种记忆}

长期对话包含时间、人物、地点、数量及其从属关系。例如，一次旅行总里程为1200英里，其中第一天的300英里已包含在总数中；若另三次旅行合计1800英里，四次旅行应为3000英里，而非3300英里。错误不仅可能来自遗漏证据，也可能来自取回后对“整体与部分”关系的误判。

论文区分事实记忆 $M$ 与元记忆 $E$：前者保存具体经历与知识，后者保存\textbf{可迁移的知识使用经验}，例如在计数前检查是否重复包含。设问题为 $q$，主方法可抽象为
\begin{equation}
\hat y=\mathrm{LLM}\bigl(\operatorname{Prompt}(q,M,E)\bigr).
\end{equation}
普通记忆问答主要提供$(q,M)$，MetaMem进一步将$E$写入提示词。优化对象是自然语言规则，基础语言模型的权重保持不变。

\subsection{从反馈到规则的自反思优化}

方法包含规则学习与冻结后的问答两个阶段。训练时，给定带标准答案的样本$(q,a,M)$，模型在当前元记忆下多次生成回答，再利用标准答案评价。反思模块归纳成功与失败原因，提出规则的添加、删除或修改，并合并批内重复与冲突的提案。验证集用于选择规则版本，测试时使用固定规则和独立问题。

\begin{center}
\small
训练数据与候选回答 $\longrightarrow$ 正误反馈与反思 $\longrightarrow$ 规则编辑与合并
$\longrightarrow$ 冻结后的元记忆 $E$ $\longrightarrow$ 新问题的事实利用
\end{center}

规则学习的主要成本来自多次回答采样、反馈、反思与编辑；冻结后的问答阶段读取事实记忆和规则。

\subsection{实验设置与结果}

主要数据集LongMemEval包含500道问题，采用五折划分，每折350题训练、50题验证、100题测试。生成模型为Qwen3-30B-A3B-Instruct与Llama3.1-70B-Instruct，事实记忆基座主要为LightMem，默认评价模型为Qwen3-235B-A22B。规则学习默认运行5个epoch，每批50题，每题采样5次。

\begin{table}[!htbp]
\centering\small
\caption{LongMemEval平均准确率（\%）}
\begin{tabular}{lrrr}
\toprule
生成模型 & LightMem & MetaMem & 增量（百分点）\\
\midrule
Qwen3-30B-A3B-Instruct & 67.50 & 71.90 & +4.40\\
Llama3.1-70B-Instruct & 66.17 & 69.08 & +2.91\\
\bottomrule
\end{tabular}
\end{table}

两种生成模型的平均准确率分别提高4.40和2.91个百分点。Qwen配置下，完整方法为71.90\%，去除反思后为70.02\%，去除规则演化后为68.33\%，说明两个环节均有贡献。

\subsection{方法局限与课设启发}

元记忆将事实召回与事实利用分开处理。自然语言规则便于检查、合并与迁移，可用于分析时间推理、计数和跨事件引用中的重复错误。

规则学习依赖带答案样本与评价模型，错误评判可能污染更新；多次采样和编辑增加优化成本。规则过度泛化、彼此冲突或积累过长也会影响回答。方法发挥作用的前提是关键事实已进入上下文，其在小模型、小样本和其他任务上的适用性仍需检验。

StructMem侧重事实与关系的组织，MetaMem侧重检索结果的使用。课程设计可据此将失败题分为证据缺失与证据利用错误，并在固定证据的条件下检查重复计数、时间误读等问题。这样的诊断有助于判断下一步应调整检索还是回答策略。

% 原文定位：原文PDF第1页，图1；原文PDF第3页，式(1)--(3)；原文PDF第5、11页，\S4与附录A.3、A.5；原文PDF第7页，表2；原文PDF第9、11--12页及本组既有代码核查。

\section{Text2Mem：记忆操作的统一表达与状态验证}
\label{sec:text2mem}

Text2Mem\upcite{wang-etal-2026-text2mem}面向自然语言驱动的记忆控制，提出统一操作语言及其基准。用户要求更新、合并或调整记忆时，系统通过结构化规范表达目标范围与参数，并检查执行后的状态是否符合预期。

\subsection{研究问题与适用范围}

长期交互不仅需要检索历史信息，还需要控制记忆的优先级、有效期、合并和删除。例如，“暂时不要再提午餐”可能意味着降低显著性，也可能被误解为删除记录。如果动作、范围与时间未被明确表达，系统即使调用成功，也可能改变错误的记忆。论文因此在自然语言与具体后端之间引入结构化操作规范。

中间规范用记忆生命周期操作表达意图，再由后端实现其语义，为不同后端提供共同的操作接口。论文主实验采用SQL/SQLite参考环境。

\subsection{操作集合、Schema与执行流程}

\paragraph{十二个操作。}
论文将操作分为编码、存储和检索三阶段。编码阶段仅有Encode；存储阶段包括Update、Label、Promote、Demote、Merge、Split、Delete、Lock、Expire九个操作；检索阶段包含Retrieve、Summarize。Encode组织内容、语义、来源与治理属性，向量生成可以延迟。Promote与Demote调整显著性，Demote可以归档而不删除；Delete处理软/硬删除，Expire设置到期规则，三者具有不同语义。

\paragraph{五字段契约。}
操作由\texttt{stage}、\texttt{op}、\texttt{target}、\texttt{args}、\texttt{meta}描述，分别对应阶段、动作、目标、专属参数和执行上下文。不同操作的必填条件不同，例如Encode创建新记忆时不必先选择已有目标。对于需要定位已有条目的操作，\texttt{target}在\texttt{ids}、\texttt{filter}、\texttt{search}、\texttt{all}四种选择方式中择一；跨字段约束限制锁定、到期和大范围写入等行为。

\paragraph{Validator--Parser--Adapter。}
Validator检查结构完整性与跨字段约束；Parser把合法实例转换为类型化操作，归一化时间和优先级等参数；Adapter再映射为具体后端动作。执行返回\texttt{ExecutionResult}，记录受影响条目和状态变化。多步骤请求可表示为保留依赖顺序的\texttt{schema\_list}。该流程检查结构和治理约束；目标选择与意图理解则由执行后的状态检查评估。

\paragraph{教学案例：提高课设进度记忆的优先级。}
以请求“提高本周课设进度笔记的重要性，其他记录不变”为例，设目标条目为17，可用STO阶段的Promote，将\texttt{ids=[17]}的\texttt{weight}设为0.9。状态检查同时覆盖目标权重、原有内容和非目标条目，确认操作只改变了指定范围。

\subsection{基准、指标与主要结果}

基准以自然语言指令、前置操作建立的初态、参考Schema和SQL状态断言组成评测实例，覆盖中英文、直接与间接指令、单操作与工作流。模型包括GPT-4o、Qwen2.5-72B-Instruct、Claude-4-Sonnet、Gemini-2.5-Pro与GPT-5，使用相同自然语言指令与提示模板。

评测包含三个指标：SMA用编辑距离或余弦相似度衡量生成Schema与参考的相似程度；ESR统计通过验证且无运行错误地执行的实例比例；EMR统计满足的SQL状态断言占全部期望断言的比例：
\begin{equation}
    \mathrm{EMR}=\frac{\sum_i\sum_{a\in A_i}\mathbf{1}[a\text{满足}]}{\sum_i|A_i|}.
\end{equation}
其中$A_i$为第$i$个实例的断言集合。EMR以断言为统计单位，断言较多的实例对总体指标影响更大。

\begin{table}[!htbp]
    \centering
    \caption{Text2Mem总体评测指标}
    \label{tab:text2mem-main}
    \small
    \begin{tabular}{lrrrr}
        \toprule
        模型 & SMA$_{\rm lev}$ & SMA$_{\rm cos}$ & ESR & EMR \\
        \midrule
        GPT-4o & 0.6776 & 0.8690 & 0.9484 & 0.2532 \\
        Qwen2.5-72B-Instruct & 0.6195 & 0.8349 & 0.9510 & 0.2535 \\
        Claude-4-Sonnet & 0.6230 & 0.8475 & 0.9594 & 0.3132 \\
        Gemini-2.5-Pro & 0.6491 & 0.8560 & 0.9479 & 0.2897 \\
        GPT-5 & 0.6115 & 0.8396 & 0.9662 & 0.2826 \\
        \bottomrule
    \end{tabular}
\end{table}

表\ref{tab:text2mem-main}列出各模型的总体指标。ESR约为94.79\%--96.62\%，EMR约为25.32\%--31.32\%，显示操作普遍可以执行，但状态效果与预期仍有较大差距。最高ESR与最高EMR分别对应GPT-5和Claude-4-Sonnet。

复杂存储和工作流的表现较弱。五个模型的Split EMR均为0；Merge仅GPT-4o为0.0465，其余为0。GPT-5的单操作EMR为0.2930，工作流EMR为0.1190，而工作流ESR为1.0000。这表明多步骤操作中的状态正确性是该基准的主要难点。

现有实验主要比较模型、操作类型和任务分组。Validator、Parser与Adapter的独立作用，以及指令复杂度对结果的影响，还需要组件消融和受控实验分析。

\subsection{方法局限与课设启发}

统一规范将动作、范围和生命周期参数显式化，便于执行前检查与执行后追踪。将结构相似、执行状态和实际效果分别计量，有助于定位目标选择、操作规划和后端执行中的问题。

方法的局限包括复杂记忆语义缺少形式保证、真实后端覆盖有限，以及尚未评测端到端长期智能体表现。状态断言的覆盖也会影响EMR：只检查目标权重，可能遗漏非目标条目被误改的情况。实际应用需要同时考虑后端适配、长期效果和错误传播。

StructMem关注事件绑定、整合与问答检索，Text2Mem关注操作表达和状态核验。与Semantic XPath的目标定位、MetaMem的事实利用结合比较，可以将记忆系统的错误依次定位到访问、推理和执行环节。

课程设计可借鉴显式输入约束与效果检查相结合的做法。若开展记忆更新实验，可固定初态、请求和模型，分别记录格式合法性、执行成功率、断言满足率与整任务通过率，并检查非目标记录是否保持完整。

% 原文定位：原文第1--3页，图1；第5、7页，\S3.3、\S4.3；第4页表1，第11页附录A；第4--5页，\S3.2；第3页图2、第5--7页\S3.3--4.1及图3；第12页附录B.1.1；后两项为本文的检查建议；第6--8页，\S4；第7--8页，\S4.1、\S4.3；第8、15页；第9页Limitations。


\section{跨论文的系统性比较与讨论}
\label{sec:synthesis}

六篇论文共同面对的问题，是如何把持续增长的历史转化为后续任务能够维护、找到并正确使用的记忆。但它们干预的环节不同：有的改变信息组织，有的改变访问范围，有的学习证据使用经验，还有的约束记忆操作的执行。因而，综合分析需要同时回答“改进了什么环节”“依赖什么条件”和“仍会在哪一步失败”，再讨论不同机制能否互补。

\subsection{统一比较框架：研究对象、机制与边界}

表\ref{tab:memory-comparison}先对齐研究对象和机制，不将六项工作视为同一个榜单上的六个可直接替换算法。表中列出各论文的主要着力点，并不表示它只具有这一项能力。

\begin{table}[htbp]
    \centering
    \caption{六篇论文的研究定位与比较边界}
    \label{tab:memory-comparison}
    \small
    \begin{tabularx}{\linewidth}{@{}p{0.16\linewidth}XXX@{}}
        \toprule
        工作 & 研究对象与表示 & 主要着力点 & 比较时的边界 \\
        \midrule
        StructMem\upcite{xu-etal-2026-structmem} & 长对话中的事件条目与跨事件综合记忆 & 保留事件内绑定，连接跨事件信息 & 综合阶段的事件重建不等于问答阶段必然采用相同扩展 \\
        Hindsight\upcite{latimer-etal-2026-hindsight} & 区分事实、经历、观察与观点的时序实体记忆 & 写入、融合检索与利用记忆推理 & 端到端系统同时包含组织、检索和推理因素 \\
        Semantic XPath\upcite{liu-etal-2026-semantic} & 树状记忆及其子树与节点 & 将结构范围约束与语义匹配结合 & 依赖可用的层级结构及正确的查询生成 \\
        HyMEM\upcite{zhu-etal-2026-hybrid} & 电脑操作中的离散经验和连续多模态表示 & 混合表示、经验更新和工作记忆刷新 & 网页交互、视觉输入与训练条件不同于纯文本问答 \\
        MetaMem\upcite{xin-etal-2026-metamem} & 事实之外的知识利用经验 & 通过反馈反思优化可复用的符号经验 & 需要区分优化阶段反馈与测试时可用信息 \\
        Text2Mem\upcite{wang-etal-2026-text2mem} & 记忆操作的结构化语言与执行状态 & 明确操作语义，验证规划和执行 & 操作正确性不能直接换算成长对话回答准确率 \\
        \bottomrule
    \end{tabularx}
\end{table}

上述定位分别依据 StructMem 第 2--3 页图 2 与公式 1--6、Hindsight 第 2--5 页图 1、Semantic XPath 第 2--5 页第 3 节、HyMEM 第 3--5 页图 1、MetaMem 第 3--5 页图 2，以及 Text2Mem 第 3--5 页图 2 和表 1。此处页码均按仓库 PDF 从第 1 页起计。

\subsection{机制之间的联系：互补发生在什么条件下}

\textbf{组织结构与访问范围。} StructMem 强调同一事件内的事实与关系绑定，以及跨事件的综合；Semantic XPath 强调查询应落在哪个结构范围。二者可以用来区分“存储时丢失关系”和“读取时选错范围”这两类问题：即使存储正确，也可能检索错会话或错节点；即使检索范围正确，也可能因为原始抽取已遗漏关系而无法回答。这里的互补是分析思路，不能据此认定把 XPath 接到事件记忆上就一定有效；还需定义结构映射、证据来源与额外查询成本。

\textbf{取回证据与使用证据。} Hindsight 的 recall 与 reflect 分别涉及检索和利用记忆推理；MetaMem 则把可复用的知识利用经验作为优化对象。它们共同提示：检索命中后仍可能出现重复计数、时间关系误读或不当综合。分析回答错误时，需要区分“证据没有进入上下文”和“证据已进入但被用错”。改善前者的检索策略与改善后者的利用经验可能互补，但效果必须分别验证，不能只观察最终分数便认定瓶颈已经解决。

\textbf{更新策略与操作语义。} HyMEM 在具体任务中决定如何合并、替换和刷新经验；Text2Mem 则强调操作应具有明确的目标、参数与可检查的执行结果。前者帮助讨论“何时更新、保留什么”，后者帮助讨论“实际执行了什么、状态是否符合意图”。但二者的表示与接口并不天然一致，尤其不能把多模态连续记忆直接当作参考 SQL 后端的一条文本记录。

\textbf{事实、观点与利用经验。} Hindsight 对事实与观点的区分，以及 MetaMem 对知识利用经验的建模，说明记忆条目的角色需要明确。世界事实、对主体的推断、操作经验和回答策略不宜无差别地作为同一类证据。它们分别在什么条件下可用于回答、何时需要更新或撤销，是比较系统设计时应追问的问题，而非仅看数据结构叫“图”还是“树”。

\subsection{共同取舍：结构化获得了什么，又付出了什么}

\textbf{关联完整性与抽取误差。} 事件、实体关系和层级有助于保留上下文，但结构本身也由抽取或整理得到。实体合并错误、时间绑定错误及摘要遗漏会进入后续检索和推理。因此应同时检查结构是否找到了来源，以及来源中的关键事实是否仍被保留，不能用“有来源 ID”替代内容正确性。

\textbf{检索覆盖与上下文开销。} 多路召回、事件补齐和更大范围的结构访问可能增加有效证据，也会引入冗余和无关信息。评价局部改动时，应在相同回答上下文预算下观察证据覆盖、回答质量及成本，避免把“给模型更多文本”的收益全部归因于结构设计。

\textbf{动态适应与错误固化。} 综合记忆、观点更新和反思经验能够适应后续任务，也可能保留错误概括或过拟合开发样本。讨论应追溯更新依赖的证据和反馈，区分系统运行中自然获得的信息与离线优化时使用的标准答案；后者不能出现在独立测试的记忆构建或检索条件中。

\textbf{表达能力与验证范围。} 明确操作 schema 能减少歧义、支持执行检查，但可解析不代表目标选择正确，可执行也不代表执行后满足用户意图。同样，能组织复杂多模态经验，不代表在更简单的文本任务上具有成本优势。需要分别报告表示、访问、推理和执行各环节的约束。

以上是用于比较的分析维度，不是声称所有论文都已经验证了这些取舍的同方向收益。具体优缺点需回到各篇论文的实验、消融和局限，避免把设计动机当成已证明的结论。

\subsection{怎样比较实验结果，才支持方法判断}

六篇论文的任务覆盖长对话问答、结构化任务状态、电脑操作和记忆指令执行。回答准确率、网页任务成功率、操作执行率和状态匹配指标测量的对象不同，不能汇总为一个“总体性能”分数。即使都报告 LoCoMo，模型、数据处理、题型范围、检索与回答预算、评分提示词等仍可能不同。

比较时应记录任务与数据划分、模型角色、是否训练或使用优化反馈、检索范围和预算、评价指标、主要成本六类条件。原论文的跨系统结果说明其各自实验条件下的表现；要判断某个结构或机制的作用，优先看该论文内部控制变量的消融。若本组复现，则所有对照共用明确的数据划分和评价口径，并说明与原论文配置的差异。

成本也要按机制分开：写入与综合开销、检索和回答开销、优化与裁判开销不是同一项。无需训练基座模型，并不等于无需多次采样与反馈；一次构建的成本，也不能直接与一次查询的成本相比。

\subsection{用共同案例检查六种研究视角}

可围绕一个明确标注为自拟的讨论案例展开：历史记录中，某人先计划周五到 A 地，后来改为周六到 B 地；随后提到从 B 地继续出发。查询要求回答最终出行时间和后续路线。

这个案例首先涉及事件内人物、时间和地点是否绑定正确，以及跨事件关系是否能被找回；其次涉及检索是否选对会话、时间和状态版本；还涉及已取回新旧计划后，是否能正确处理变更，而非把两个版本拼成矛盾答案。若用户再要求暂时不提及这次出行，则还要明确这是降低检索优先级、临时隐藏还是删除，操作范围和期限分别是什么。

讨论时应由各论文负责人说明：其方法直接处理哪一部分、依赖什么输入、哪一部分并未给出保证，并指出原文依据。这个案例用于检验解释是否一致，不是论文原例，也不是已运行的六系统对照实验；尤其不能把纯文本例子直接视为验证了 HyMEM 的多模态机制。

\subsection{从比较形成课设选择与团队结论}

综合分析最终应落实为有依据的选择：本组要优先解决信息组织、检索范围、证据利用还是操作执行中的哪一个问题；为何选择某项工作作为主复现对象；哪些机制只作为参考，哪些因任务或资源不匹配而暂不纳入。选择某篇作为主线并不意味着其他五篇只作背景罗列，还应说明它们如何帮助定义错误类型、评价指标或可选改进。

此部分等待实际团队讨论后补充：保留有原文或实验依据的共同结论；对尚有分歧的问题写明分歧点与需要核查的证据。若提出机制组合，应分别列出结构适配、信息边界和对照实验，而非直接宣称已提出有效的新方法。最终报告与 PPT 应呈现这些经过讨论的取舍，并在成员分工中记录实际参与的比较、核对和整合工作。


\clearpage
\section{总结}
六篇论文把长期记忆的困难分解到了不同位置。StructMem通过事件绑定与跨事件综合保留关联，Hindsight区分记忆角色并组织多路召回，Semantic XPath将访问约束落实到层级范围，HyMEM结合多模态经验与两层更新，MetaMem学习事实利用规则，Text2Mem则把操作表达与状态核验分开。它们共同说明，记忆系统需要从来源、组织、访问一直追踪到使用或执行结果；一个环节的改进不能自动替代其他环节的正确性。

从证据上看，结构化表示能够支持更明确的访问和诊断，但也引入抽取、维护、检索扩展或优化反馈的成本。原文的总体提升不意味着所有题型和配置均获益，较少的输入tokens也不等同于端到端成本更低。可靠比较应先对齐任务、模型、反馈、预算和评价口径，再结合受控消融与具体失败案例解释差异。

据此，本组延续以StructMem进行小规模方法复现的既定范围，将其余五篇作为检索范围、时间更新、证据利用及操作核验的参照。阅读报告给出的进一步判断是：先定位错误发生在哪一层，再决定是否引入相应机制，并保留验证该机制所需的对照条件。这一判断为后续实验提供分析框架，不预设机制组合必然有效，也不把个人试点或作者结果合并为本组成绩。

\bibliographystyle{jos}
\bibliography{references}

\section*{成员分工与贡献}
待填写：六位成员姓名、负责论文、报告章节、PPT 与讲解工作及实际贡献百分比，总和为 100\%。
\end{document}
